% Standalone insert. Requires amsmath, amssymb, amsthm; see insert_driver.tex.
% Namespaced labels/macros deliberately avoid the canonical manuscript's labels.
\providecommand{\NLFP}{\operatorname{FP}}
\providecommand{\NLLi}{\operatorname{Li}}
\providecommand{\NLA}{\mathcal A}
\providecommand{\NLT}{\mathcal T}
\newtheorem{nonlinearthm}{Theorem}[section]

\section{Nonlinear coordinate transport of Stieltjes finite parts}
\label{nonlinear:sec:transport}
All finite parts in this section use a cutoff in the named coordinate and
logarithmic scale one. Thus $\NLFP\int_0^1x^{-1}\log^n x\,dx=0$, while for
$k\ge2$ its value with $x^{-k}$ is $-n!/(k-1)^{n+1}$.
A density pullback is defined by
$\langle\phi^*T,h\rangle=\langle T,h\circ\phi^{-1}\rangle$;
its ordinary coefficient is $\phi'f(\phi)$.

\begin{nonlinearthm}[Complete local coordinate correction]
\label{nonlinear:thm:local}
Let $\phi$ be a smooth increasing endpoint-preserving diffeomorphism of
$[0,1]$, with $\lambda=\phi'(0)>0$. Put $r(y)=\phi(y)/y$ and
$L(y)=\log r(y)$. For $k\ge1$, $n\ge0$,
\begin{align*}
\NLA_\phi^{k,n}
&:=\phi^*\NLFP_x(x^{-k}\log^n x\,dx)
-\NLFP_y(\phi'\phi^{-k}\log^n\phi\,dy),\\
\langle\NLA_\phi^{k,n},h\rangle
&=\frac1{n+1}[y^{k-1}]\phi'(y)r(y)^{-k}L(y)^{n+1}h(y).
\end{align*}
In particular the coefficient of $\delta^{(j)}$ is
\[
\frac{(-1)^j}{(n+1)j!}[y^{k-1-j}](\phi'r^{-k}L^{n+1}),
\qquad 0\le j\le k-1.
\]
The formula extends linearly to finite logarithmic singular expansions
modulo locally integrable remainders. Its discrepancy obeys
\[
\NLA_{\phi\circ\psi}(f\,dx)
=\psi^*\NLA_\phi(f\,dx)+\NLA_\psi(\phi^*(f\,dx)).
\]
\end{nonlinearthm}
\begin{proof}
The meromorphic distribution $y^{-k-u}$ has residue
$-h^{(k-1)}(0)/(k-1)!$ on a test function $h$. Taylor subtraction through
degree $k-1$ leaves an integral holomorphic for $|\Re u|<1$; each fixed
spectral derivative is dominated by an integrable power times finitely
many logarithms. Its regular coefficients are
$(-1)^n\NLFP(y^{-k}\log^n y)/n!$.
Under $x=\phi(y)$ the full numerator is
$A(u,y)=\phi'r^{-k}e^{-uL}$.
The regular coefficients give the coordinate finite part of the composed
logarithmic singularity. The additional coefficient comes from
$-[y^{k-1}](A(u,y)h(y))/u$ and is
$[y^{k-1}](\phi'r^{-k}L^{n+1}h)/(n+1)$ after multiplying by $(-1)^nn!$.
This proves the action formula and the delta coefficients. Integrable
remainders obey ordinary substitution. The composition rule follows by
adding and subtracting $\psi^*\NLFP(\phi^*(f\,dx))$.
\end{proof}

An existing comparison $K=\NLFP(f\,dx)+\kappa\delta^{(d)}$ consequently
transports to
\[
\phi^*K=\NLFP(\phi'f(\phi)\,dy)
+\NLA_\phi(f\,dx)+\kappa\phi^*\delta^{(d)}.
\]
For $d\ge1$ the last term has the explicit derivative expansion
\[
\phi^*\delta^{(d)}=
\sum_{j=1}^d(-1)^{d+j}\frac{(d-1)!}{(j-1)!}
[y^{d-j}]\left(\frac y{\phi(y)}\right)^d\delta^{(j)}.
\]
Indeed, writing $\tau=\phi^{-1}$, its action is
$(-1)^dd![x^d]h(\tau(x))$; Lagrange inversion yields the displayed
coefficients. Also $\phi^*\delta=\delta$. For a second endpoint use
$1-\phi(1-y)$ and reflect its delta derivatives with sign $(-1)^j$.
This is transport of the specified distribution, not a claim that nonlinear
pullback preserves circular convolution.

\subsection{Harmonic Stieltjes specialization}
Use
\[
\zeta(1+u,x)=\frac1u+\sum_{m\ge0}\frac{(-1)^m}{m!}\gamma_m(x)u^m,
\qquad \gamma_0=-\psi.
\]
Let $P_p(u)=\prod_{j=1}^p(1+u/j)=\sum_je_{p,j}u^j$, with $P_0=1$ and
$e_{p,j}=0$ outside $0\le j\le p$. Its coefficients are finite harmonic
polynomials because
$\log P_p(u)=\sum_{r\ge1}(-1)^{r-1}H_p^{(r)}u^r/r$.
Define
\[
\NLT_{m,p}(L)=m!\sum_{j=0}^{\min(m,p)}
\frac{(-1)^je_{p,j}L^{m-j+1}}{(m-j+1)!}.
\]
\begin{nonlinearthm}[All-index Stieltjes correction]
\label{nonlinear:thm:stieltjes}
For every $m,p\ge0$,
\[
\big\langle\NLA_\phi(\gamma_m^{(p)}\,dx),h\big\rangle
=(-1)^pp![y^p]\phi'r^{-p-1}\NLT_{m,p}(\log r)h.
\]
If $\phi(y)=y+ay^2+O(y^3)$ and $p\ge1$, this correction vanishes for
$m\ge2p$, while at the last possible index it equals
\[
\NLA_\phi(\gamma_{2p-1}^{(p)}\,dx)
=\frac{(2p-1)!}{p!}a^p\delta.
\]
In particular $\NLA_\phi(\psi'\,dx)=a\delta$.
\end{nonlinearthm}
\begin{proof}
The Hurwitz shift gives the exact singularity
\[
\gamma_m^{(p)}(x)=(-1)^pp!\sum_{j=0}^{\min(m,p)}
\frac{(-1)^jm!e_{p,j}}{(m-j)!}\frac{\log^{m-j}x}{x^{p+1}}
+\gamma_m^{(p)}(1+x).
\]
Apply Theorem~\ref{nonlinear:thm:local} term by term. When the slope is one,
$\log r=ay+O(y^2)$. For $m\ge2p$ the smallest logarithmic power in
$\NLT_{m,p}$ is at least $p+1$, so its $y^p$ coefficient vanishes. At
$m=2p-1$ only $j=p$ survives; use $e_{p,p}=1/p!$.
\end{proof}

\subsection{Nonlinear circle averages}
Let $\chi_\lambda:[0,1]\to[0,1]$ be the continuous increasing lift of
$\tan(\pi\chi_\lambda(y))=\lambda\tan(\pi y)$, where $\lambda>0$.
It equals $\pi^{-1}\arctan(\lambda\tan\pi y)$ below $1/2$, equals $1/2$
there, and requires adding one to that expression above $1/2$.
Put $q=(1-\lambda)/(1+\lambda)$ and $R_\lambda(y)=\chi_\lambda(y)/y$.
The inverse density is the Poisson kernel
\[
(\chi_\lambda^{-1})'(x)=\Pi_q(x)
=\frac{1-q^2}{1-2q\cos(2\pi x)+q^2}.
\]
The Taylor series of $R_\lambda$ is even.

For spectral coefficient extraction set
\[
F_p(u;q)=2(2\pi)^p\Gamma(1-u)(2\pi)^u
\cos\frac{\pi(p+u)}2\,\NLLi_{-p-u}(q).
\]
\begin{nonlinearthm}[Exact nonlinear Stieltjes averages]
\label{nonlinear:thm:circle}
For $m,p\ge0$,
\begin{align*}
\NLFP\int_0^1\gamma_m^{(p)}(\chi_\lambda(y))\,dy
={}&(-1)^{m+1}m![u^{m+1}]
\bigl(F_p(u;q)-P_p(u)F_p(0;q)\bigr)\\
&-(-1)^pp![y^p]R_\lambda^{-p-1}
\NLT_{m,p}(\log R_\lambda).
\end{align*}
The last line vanishes for odd $p$. At Stieltjes index zero,
\[
\NLFP\int_0^1\psi^{(2j+1)}(\chi_\lambda(y))\,dy
=(-1)^{j+1}\pi(2\pi)^{2j+1}\NLLi_{-2j-1}(q).
\]
These last values are rational functions of $q$ times the displayed power
of $\pi$.
\end{nonlinearthm}
\begin{proof}
The canonical periodic Stieltjes generator has nonzero Fourier coefficients
\[
(2\pi i n)^p\left[\Gamma(-u)(2\pi i n)^u+\frac{P_p(u)}u\right]
\]
and zero mean. This follows from the Hurwitz Fourier formula and full
endpoint-residue subtraction; it is the previously established
fixed-coordinate identity. Pairing it with $\Pi_q$ gives
$-(F_p(u;q)-P_p(u)F_p(0;q))/u$. Exponential Fourier decay justifies every
fixed spectral derivative. To change to the unweighted composition
average, apply Theorem~\ref{nonlinear:thm:stieltjes} with
$\phi=\chi_\lambda$ and test function $1/\chi_\lambda'$.
This cancels the Jacobian in the contact coefficient. Evenness of
$R_\lambda$ gives the odd-order zero; differentiating the spectral factor
at $u=0$ and using $\gamma_0=-\psi$ proves the polygamma formula.
\end{proof}

Ordinary convergent specializations, with
$\NLLi_s^{[1]}(q)=\partial_s\NLLi_s(q)$, are
\begin{align*}
\int_0^1\left[\psi(\chi_2(y))+\frac1{2y}\right]dy
&=-2\NLLi_0^{[1]}(-1/3)-\frac{\gamma+\log\pi}{2},\\
\int_0^1\left[\psi'(\chi_2(y))-\frac1{4y^2}\right]dy
&=\frac14+\frac{3\pi^2}{8},\\
\int_0^1\log\Gamma(\chi_\lambda(y))\,dy
&=\frac12\log\frac{\pi(\lambda+1)}\lambda.
\end{align*}
The first two follow by subtracting their complete singular heads and using
the monomial finite parts stated above. The last follows independently
from Gamma reflection, symmetry of the map and the Poisson average of
$\log\sin(\pi x)$. The full article supplies the all-index ordinary
subtraction procedure, normalized antiderivative hierarchy, source
attribution and verification artifacts. These statements do not resolve
unrestricted multiple collision restrictions or the remaining Gaussian
$S_6$ reduction.
